\section{Realización General del Proyecto}
    Las implementaciones de las transformaciones afines se realizaron en NumPy \parencite{numpy_library}, 
    las cuales fueron verificadas de forma visual sobre el conjunto de imágenes empleado 
    para este proyecto. Esta verificación permitió mejorar la implementación de forma progresiva 
    para ajustarlas al resultado de la transformación esperada de forma individual, principalmente central el 
    efecto de la imagen respecto a su origen. De esta verificación se determinó calcular primero 
    la matriz de la transformación afín resultante de aplicar varias transformaciones de forma 
    consecutiva con el fin de generar un mejor resultado con una menor pérdida de información por 
    una transformación incompleta.\\

    Para aplicar correctamente la transformada afín, se empleó de la transformada inversa 
    junto con la interpolación bilineal para estimar la intensidad del color en el pixel 
    de la imagen resultante (transformada), este procedimiento se basa en lo visto en clase \parencite{tinoco2026pdi}, 
    siguiendo esta metodología se lograron los resultados técnicos esperados de una imagen 
    transformada con un buen tiempo de ejecución, esto al considerar que todo las operaciones 
    pesadas son ejecutadas como multiplicación de matrices en NumPy y solamente existe un 
    ciclo \textit{for} de Python para evaluar cada pixel de la imagen resultante.\\ 